\clearpage\thispagestyle{empty}
\pdfbookmark[0]{Quellen zu Pallas (Brief vom 4. August 1812)}{pallas1812aug4-sources}
\begin{center}\textbf{Quellen zu Pallas (Brief vom 4. August 1812)}\end{center}
\noindent Diese KI-Übertragung ins heutige Deutsch folgt dem Brief von Gauß vom 4. August 1812 in den \href{https://gdz.sub.uni-goettingen.de/id/PPN236007467}{\textit{Werken}, Band~VI} (Göttingen, 1874), S.~354--356. Der Brief erschien zuerst in der \textit{Monatlichen Correspondenz zur Beförderung der Erd- und Himmelskunde}, Band~XXVI, August 1812, Nr.~XIII, S.~\href{https://api.digitale-sammlungen.de/iiif/image/v2/bsb10538615_00207/full/1600,/0/default.jpg}{199}, \href{https://api.digitale-sammlungen.de/iiif/image/v2/bsb10538615_00208/full/1600,/0/default.jpg}{200}, \href{https://api.digitale-sammlungen.de/iiif/image/v2/bsb10538615_00209/full/1600,/0/default.jpg}{201}, \href{https://api.digitale-sammlungen.de/iiif/image/v2/bsb10538615_00210/full/1600,/0/default.jpg}{202}, \href{https://api.digitale-sammlungen.de/iiif/image/v2/bsb10538615_00211/full/1600,/0/default.jpg}{203}. Das digitalisierte Exemplar stammt aus der Bayerischen Staatsbibliothek.
\par\smallskip\noindent Für die wahre geozentrische Breite bei der Opposition stehen in der Zeitschrift $48\;16\;26.1\;\mathrm{N}.$ und in der Werkausgabe $48\;18\;26.1\;\mathrm{N}.$ Die Werkausgabe enthält außerdem vier Punkte hinter der Deklination vom 9. April, zwei waagerechte Striche vor der ersten Deklinationsdifferenz sowie die unterschiedlichen Schreibungen Burkhardt im Text und Burckhardt in der Tabelle. Die Übertragung folgt dem Brief in der Werkausgabe. KI-Übertragung und KI-Satz.
\clearpage
